\section{Giriş}
\indent\indent Bu çalışma, Tanzimat Dönemi ile birlikte hız kazanan Batı tarzı müesseselleşme sürecinin Osmanlı ceza hukuku üzerindeki etkilerini incelemeyi amaçlamaktadır. Tanzimat'tan sonraki yirmi yıllık
süreçte 1840 Kanunnâme-i Hümâyunu, 1851 Kanun-ı Cedîd ve 1858 Kanunnâme-i Hümâyunu olmak üzere üç ceza kanunnâmesi kaleme alınmıştır. Bu süreçte, ceza hukuku şer`î hukuktan neredeyse tamamen 
ayrılmıştır. Şer`i hukuk, ceza hukuku üzerine net bir kodifikasyon ihtiva etmediği için verilen kararlarda kadının yorumu ne önemli unsurdu. Bu karara itirazı yönetecek temyiz mekanizması ise hemen hemen yoktu.
Bu süreç ile ceza hukuku, kazuistikkuram üzerine oturtularak yeni mahkemeler teşkil edilmiş ve temyiz mekanizması oluşturulmuştur. Çalışma, bu değişim sürecinin nedenlerini ve meydana getirdiği bürokrasi çatışmalarını, mezkur kanunnâmeler üzerinden incelemektedir.

\subsection{Tanzimat Öncesi Ceza Hukuku}

\indent\indent Ceza hukukunu inclemeden önce, genel hatlarıyla Osmanlı Devleti'nin hukuk sistemine göz atmak faydalı olacaktır. Osmanlı hukuku, şer`î hukuk ve örfî hukuk olmak üzere iki ana kaynaktan beslenmekteydi. 
Şer`î hukuk kaynağını öncelikli olarak kitap ve sünetten almaktadır. Kitap ve hadisin eksik ya da muğlak bıraktığı kısımlarda ise icmâ, kıyas, istihsân, maslahat gibi yöntemler kullanılmıştır. Osmanlı Devleti
şer`î hukukta ağırlıklı olarak Hanefi mezhebini benimsemiştir. Örfî hukukun kaynağını ise siyasi otorite, yani mahalli örf, geçmiş uygulamalar ve en önemlisi olarak padişah iradesi teşkil etmektedir.\footfullcite[90]{osmanli_hukuku}
Şer`î hukuk ve örfî hukuk birbiriyle çatışan iki farklı uygulama olarak anlaşılmamalıdır. Bu iki hukuk sistemi birbirini tamamlayan ve dengeleyen unsurlar olarak görülmelidir. Şer`î hukuk ağırlıklı olarak şahıs, aile, miras, eşya, borçlar,
ticaret ve yargılama gibi alanlarda düzenlemeler yaparken, kısmen de ceza, arazi ve vergi hukuku alanlarına müdahil olmuştur.\footcite[95]{osmanli_hukuku} Örfî hukuk ise şer`î hukukun düzenlemediği ya da kısmen boş bıraktığı alanlarda doldurmaya çalışmıştır. Örfî hukuk,
devletin siyasi ve idaari yapısı, vergi, arazi ve ceza hukuku alanlarında düzenlemeler getirmiştir.\footcite[97]{osmanli_hukuku} Buradan da anlaşılacağı üzere, ceza hukuku iki sistem tarafından da düzenlemelere maruz kaldığı için
en çok tartışılan alanların başında gelmiştir.

Tanzimat öncesi Osmanlı ceza hukukunun işleyişini, devletin teb`aları açısından incelemek faydalı olacaktır. Çok büyük farklılıklar olmamakla birlikte, ceza hukukun uygulanışı müslümanlara, zımmîlere ve müste'menlere karşı ufak
değişiklikler göstermekteydi.
\subsection{Müslümanlara Uygulanan Ceza Hukuku}
\indent\indent Tanzimat öncesi Osmanlı ceza hukukunda müslümanlar şer`î hukuk ve örfî hukuk birlikte kullanılmaktaydı. Bu iki hukukun ortaklaşa işleyişini
anlamak için suçlara ve cezalara birlikte bakmak faydalı olacaktır. Şer`î hukukta suçlar genel olarak had ve tâzir olmak üzere ikiye ayrılmaktaydı. Had,
adam öldürme ve vücut bütünlüğüne yönelik işlenen suçlardı. Bunların dışında kalan suçlar ise tâzir kapsamında değerlendirilmekteydi.\footcite[99]{osmanli_hukuku}

Şer`î hukukta had suçlarının cezaları kısastır. Örfî hukuk ise had cezalarına uygulanacak para cezlarını belirlemiştir. Fatih Kannunamesi'nde bu durumla alakalı olarak şu ifade yer almaktadır:
\begin{quote}
Eğer adam öldürse yerine kısas etmeseler kan cürmi bay olup bin akçeye dahi ziyadeye gücü yeterse dört yüz, eğer altı yüze gücü yeterse iki yüz akçe, ondan aşağı halli olursa yüz akçe ve fakir olursa elli akçe alına.\footcite[101]{osmanli_hukuku}
\end{quote}

Hırsızlık suçunda ise kısas gerektiren hırsızlıklarda şer`î hukuk temel alınırken, para cezası gerektiren hırsızlıkların cezası örfî hukuk tarafından belirleniyordu. Bu noktada hırsızlık suçlarına uygulanan 
cezaların, hırsızlığın nitekli ya da basit olmasına göre değiştiği görülmektedir. Şer`î hukukta hırsızlığın nitekli olup olmaması, miktarına göre belirlenmiştir. Hanefî hukukuna göre bu miktar 10 dirhem gümüştür.
Kānunnâme-i Sultan Selim Han'da şöyle bir hükme rastlıyoruz:
\begin{quote}
Eğer bir kimesne kaz ya tavuk ya ördek uğurlasa kadı tâ'zir edüb iki ağaca bir akçe cürüm alına. Eğer bir kimesne kovan ya koyun ya kuzu uğurlasa nisabına yetişmese yine tâ'zir edüb ağaç başına bir akçe cürüm alına. Eğer bir kişi at ya katır ya merkep, uğurlasa elin keseler yahud iki yüz akçe cürüm alına.\footcite[105]{osmanli_hukuku}
\end{quote}

Bu noktada kürek cezasına değimenin de faydası olacaktır. Kürek cezası, şer`î hukukta yer almayan ancak örfî hukukun ortaya koyduğu bir ceza türüdür. Teşkilatlanmanın henüz beylik aşamasında olduğu, 14. yüzyıl 
gibi erken bir dönemde Kara Mürsel Bey'in iştiraki ile donanma faaliyetleri başlamış olsada, gerçek anlamda bir donanma teşkilatlanması 16.yüzyılda vuku bulacaktır. Büyüyen teşklatlanma ile de donanma için gerekli 
olan insan gücünün sağlanması gerekecektir. 16.yüzyıl itibariyle şer`î hukukun getirdiği kısas cezaları daha sık şekilde kürek cezasına çevrilmeye başlanacaktır.\footcite[110]{osmanli_hukuku}

Bu bilgiler ışığında, şer`î hukuk ve örfî hukukun birbiriyle çatışan değil, birbirini tamamlayan unsurlar olduğu açıktır. Örfî hukuk, şer`î hukukun getirdiği suç tanımlarına müdahil olmazken, cezaların tayiyininde 
devreye girmiştir. Örfî hukuk, had ve tâzir suç tanımlarını açıklığa kavuştururken, para cezalarının düzenlenmesinde rol oynamıştır. Değişen zaman ve ekonomik şartlar göz önünde bulundurulduğunda, bu düzenlemelerin 
siyasi otorite için bir tercihten çok bir zorunluluk olduğu anlaşılmaktadır.
\subsection{Zımmîlere Uygulanan Ceza Hukuku}
\indent\indent İslam dininde zorlama yoktur. Bu ilke doğrultusunda, dârülharbden ele geçirilen bölgeler sakinlerininin iki seçeneği bulunurdu. İsteyenler bölgeyi terk edip, dârülislam sınırları dışına çıkabilirlerdi. 
Kalanlar ise İslam hukukunun şartlarına uymayı ve askerlikten muaf kalmayı kabul ederek, cizye vergisi yükümlülüğünü üstlenirlerdi. Taraflar arasında yapıldığı varsayılan bu anlaşmaya zimmet anlaşması, kabul edenlere de zımmî denirdi.\footcite[293]{osmanli_hukuku}

Dârülislam sınırları içerisinde islam hukuku geçerliydi. Zımmîlerin davalarından da şer`î mahkemeler sorumluydu. Fakat Osmanlı Devleti, din ve vicdan hürriyeti kapsamında zımmîlere hukuki ve kazai muhtariyet 
tanımıştır. Bu sayede zımmîler kendi aralarındaki davaları cemaat mahkemelerinde görebiliyorlardı. Cemaat mahkemelerinin yaptırım gücü bulunmadığı için daha çok bir hakem kurulu işlevi görüyordu. Taraflar 
cemaat mahkemesinin verdiği hükmü kabul ederlerse, şer`î mahkemeler müdahil olmuyordu. Öte yandan taraflardan biri cemaat mahkemesinin kararını kabul etmezse, davayı şer`î mahkemelere taşıyabiliyordu. Şer`î 
mahkemeler, bu yapıda için itiraz mahkemesi işlevi görüyordu.

Cemaat mahkemelerine verilen imtiyazın ilk ne zaman ortaya çıktığı bilinmemektedir. Ancak bu imtiyazın sınırları zamanla daraltılmıştır. İslam hukukçuları zamanla bu imtiyazın sadece ahvâl-i şahsiyye dedikleri 
konularla sınırlı kalması gerektiği fikrini belirlemişlerdir. Bu yaklaşımın sonucu olarak, cemaat mahkemeleri sadece evlilik, boşanma, miras ve vasiyet gibi konularda yetkili kılınmıştır. Öte yandan da, cemaat 
mahkemelerinin zaman zaman ceza hukuku kapsamındaki davalara da baktığı bilinmektedir.\footcite[297]{osmanli_hukuku}

Şer`î mahkemelerde, gayrimüslimlere uygulanan ceza hukuku müslümanlara uygulanan ceza hukuku arasında bazı farklar bulunmaktaydı. Dinleri içki içmeye izin verdiği için kamu düzenini bozmadıkları sürece had cezası 
uygulanmamaktaydı. Zina suçunda ise, recm ya da had cezası yerine daha hafif bir ceza olan dayak cezası uygulanıyordu. Öte yandan domuz ve şarap gibi İslam hukukunda mütekavvim olmayan mallar, gayrimüslimler için 
mütekavvim sayılıyordu.\footcite[315]{osmanli_hukuku}

Gayrimüslimlere uygulanan bu imtiyazların, zaman içinde belirli grupların otonom bir yapı oluşturmalarına yol açtığı görülmektedir. Bu gruplara millet ismi verilmiştir. Osmanlı Devleti içindeki bütün ortodokslara
\textit{Rum milleti} denmiş ve Fener patrikhanesi bu milletin merkezi olmuştur. Benzer bir yapılanma \textit{Ermeni milleti} için de gerçekleşmiştir. Bu durum zaman içinde Osmanlı Devleti için sorun 
oluşturmaya başlamıştır. 1774 Küçük Kaynarca Anlaşması ile Osmanlı Devleti, dahilindeki ortodoksların din ve vicdan hürriyetinin korunması konusunda garanti vermiştir. Bu tarihten sonra gayrimülsimlerin durumları uluslarası bir hal almıştır.\footfullcite[148]{turk_hukuk_tarihi}
\subsection{Müste'menlere Uygulanan Ceza Hukuku}
\indent\indent İslam ülkesi vatandaşı olmayıp, devlet tarafından verilen geçici süre izinle İslam ükesinde bulunan kişiler müste'men olarak adlandırılıyorlardı. 
Bu kimselerin İslam ve Osmanlı hukuku içindeki statüleri büyük ölçüde zımmîlere benziyordu. Hanefî mezhebine göre, şahıs haklarını ihlal eden suçlarda 
müste'menlere İslam hukuku uygulanırdı. Bu bağlamda müslüman ya da zımmîyi öldüren müste'mene, şartlara bağlı olarak kısas cezası uygulanırdı. Zina ve 
içki içme gibi suçlarda ise İslam hukuku uygulanmazdı. Öte yenadan Şâfiî, Mâlikî ve Hanbelî mezhepleri Hanefî  mezhebinden farklı olarak, her türlü suçta 
müste'menlere İslam hukukunun uygulanmasını savunuyorlardı.\footcite[324]{osmanli_hukuku}

16.yüzyılda uygulamaya konulan kapitülasyonlar ile yabancı devletler hukuki alanda da bazı imtiyazlar elde etmişlerdir. Fransa'ya 1535 yılında verilen 
kapitülasyonlar ile Fransa devletine tabiî tüccarların kendi aralarındaki anlaşmazlıklarda, bu devletin konsoloslukları yetkili kılınmıştır. Kapitülasyonlar,
zaman içinde diğer devletlerin de faydalanacağı şekilde genişletilmiştir. Böylece, müste'menler giderek Osmanlı mahkemelerinin yargı alanının dışında kalmaya 
başlamışlardır. Müste'menlerin Osmanlı teb`asına karşı işlediği suçlarda Osmanlı mahkemeleri yetkili olmaya devam etse de, bazı davaların Dîvân-ı Hûmâyun'da 
görülmesi ve yargılama esnasında bir elçilik görevlisinin hazır bulunması gibi bir takım ayrıcalıklar tanınmıştır.\footcite[325]{osmanli_hukuku}